% !TEX root = ../main.tex
\section{Propuesta de un archivo indexado que permita acelerar las búsquedas}

Recorrer secuencialmente un archivo de cientos de miles de registros para
ubicar un solo estudiante cuesta $O(n)$: es aceptable con 25\,000
registros, pero se vuelve lento a la escala que el caso de estudio
anticipa. La solución no es cargar todos los registros a memoria (la
sección 7 ya descarta esa opción por el espacio que ocuparía), sino
mantener en RAM una estructura mucho más liviana que solo indique dónde
está cada registro en disco: un archivo indexado.

\subsection{Funcionamiento general de un índice ISAM}

El patrón adoptado corresponde a ISAM (\emph{Indexed Sequential Access
Method}), la técnica que usaban los primeros motores de bases de datos
relacionales antes de que los árboles B y B+ se volvieran el estándar.
ISAM separa dos responsabilidades que un árbol balanceado resuelve juntas:
mantiene los datos ordenados en un archivo secuencial, y mantiene aparte
una estructura de índice, también ordenada, que guarda únicamente la
llave y la posición física del dato. Con esa estructura se puede saltar
directo al registro deseado en vez de recorrer el archivo completo.

La parte que le da a ISAM su nombre completo es cómo maneja las
inserciones. En vez de reescribir el archivo ordenado cada vez que entra
un estudiante nuevo, esas altas se acumulan en una zona aparte, sin
ordenar, y se integran a la zona ordenada más adelante, en un solo
proceso, cuando esa zona aparte ya acumuló suficientes.
% TODO: ampliar con el detalle de bloques y paginas si se decide profundizar mas
% en el funcionamiento clasico de ISAM (IBM) para la version final.

\subsection{Por qué ISAM y no un árbol}

Un árbol B o B+ resuelve el mismo problema reacomodándose automáticamente
en cada inserción, sin necesidad de un proceso de reorganización aparte.
Se descartó por dos razones concretas de este proyecto, no por evitar la
complejidad de implementarlo:

\begin{itemize}
    \item El patrón de uso del sistema no lo necesita. Un árbol
    balanceado justifica su costo de mantenimiento cuando las inserciones
    y las búsquedas están genuinamente mezcladas todo el tiempo. Aquí las
    búsquedas (matrícula, consulta de expediente, reportes) ocurren mucho
    más seguido que las altas de estudiantes nuevos, que llegan en lotes
    concentrados durante los períodos de admisión. Pagar el costo de
    reacomodar un árbol en cada inserción, para un patrón de uso que es
    mayoritariamente lectura, no se justifica.
    \item El costo que ISAM sí paga (juntar la zona ordenada con la zona
    de altas nuevas) es predecible y se puede programar para el momento en
    que menos afecta al sistema: cuando la zona de altas nuevas crece lo
    suficiente, o al apagar el sistema (ver la subsección de
    reorganización más abajo y el ADR 004). Un árbol no permite posponer
    ese costo de la misma forma, porque se reacomoda en cada operación
    individual.
\end{itemize}

\subsection{Índice primario: carné}

El índice primario ordena por carné y se divide en dos zonas, tal como se
documenta en el ADR 001:

\begin{itemize}
    \item Zona primaria: arreglo ordenado, donde se busca con búsqueda
    binaria en $O(\log n)$.
    \item Zona de altas nuevas: arreglo sin ordenar donde se agregan los
    estudiantes recién registrados en $O(1)$, revisado de principio a fin
    cuando la zona primaria no encuentra la llave buscada.
\end{itemize}

Cada entrada del índice ocupa únicamente el carné, la posición en disco y
el estado, no el registro completo. Eso es lo que permite que el índice
quepa en RAM aunque los datos completos no quepan (sección 7).

\subsection{Índices secundarios: identificación y apellido}

El caso de estudio exige buscar también por identificación y por
apellido, ninguna de las dos es la llave primaria. En vez de recorrer el
archivo completo para esas búsquedas, se agregan dos índices adicionales
(ADR 003), los dos apuntando a la misma posición física en
\texttt{identidad.dat}, sin duplicar el registro del estudiante:

\begin{itemize}
    \item Índice por identificación: funciona igual que el índice de
    carné, con zona primaria y zona de altas nuevas, porque buscar por
    identificación pasa tan seguido como buscar por carné, por ejemplo
    durante matrícula o para verificar la identidad de un estudiante.
    \item Índice por apellido: no se actualiza en cada estudiante nuevo
    como los otros dos. En vez de eso, se arma de nuevo por completo la
    primera vez que alguien pide una búsqueda o un reporte por apellido
    después de que hubo cambios desde la última vez que se armó, y
    también se arma de nuevo durante el proceso de apagado. Se hace así
    porque las búsquedas por apellido son típicamente para generar
    reportes, algo que ocurre de vez en cuando, no una operación que se
    repite en cada uso del sistema. Armar este índice de nuevo en cada
    estudiante registrado pagaría un costo de $O(n \log n)$ cada vez, por
    un beneficio que casi nunca se usa de inmediato.
\end{itemize}

Esta forma de trabajar (un solo dato físico, varios índices apuntando a
él) es el mismo patrón que un índice secundario en un motor de bases de
datos SQL: la tabla vive en un solo lugar, los índices solo indican dónde
buscar.

\subsection{Reorganización y persistencia}
\label{sec:reorganizacion}

Cuando la zona de altas nuevas del índice primario o del índice de
identificación llega al 10\% del tamaño de su zona primaria, se ejecuta
una reorganización con el sistema encendido: se juntan ambas zonas en
$O(n \log n)$ y la zona de altas nuevas queda vacía otra vez. Esta
reorganización no toca los archivos de datos en disco, solo la estructura
de índice en RAM.

Al apagar el sistema ocurre un proceso distinto, descrito en el ADR 004:
primero se purgan los registros inactivos de los tres archivos de datos,
y después el índice completo se arma de nuevo desde cero a partir del
archivo ya compactado, en vez de juntar la zona de altas nuevas que
existía antes de purgar. La razón es que la purga cambia la posición
física de todos los registros que estaban después de uno eliminado, lo
cual deja invalido cualquier índice armado antes de purgar, sin importar
si ya estaba reorganizado o no.

El índice completo (zona primaria, zona de altas nuevas y los dos índices
secundarios) se guarda en el archivo \texttt{estudiantes.idx} al apagar,
para que la próxima vez que se encienda el sistema alcance con cargar el
índice ya armado en $O(n)$, en vez de reconstruirlo desde los archivos de
datos.
